\section*{Declaration on the Use of AI Tools}
Generative artificial intelligence was used to support this thesis. 
The tool was \emph{Claude Code} (Anthropic), based on the \emph{Sonnet 4.6}, \emph{Opus 4.8}, and \emph{Opus 5.0} models, used during 2026.
It was employed for three purposes: the revision and copy-editing of prose to improve clarity, coherence, and internal numerical consistency; the drafting of sentences that were subsequently reviewed and reworked by the author; and software-engineering and data-analysis support, namely assistance in debugging the dataset-construction and evaluation pipeline, in computing the reported statistics, and in generating the result and appendix tables from experimental data.

This assistance concerns the chapters on the implementation of the theory and on type prediction with LLMs, together with the abstract, conclusion, contributions, and appendices; it does not concern the formalisation chapter, which was authored independently. 
No experimental result was generated by the tool: all training runs and evaluations were carried out by the author and recomputed directly from the data and the compiler. 
Output produced by the tool was checked against the source code and the compiler's behaviour, corrected where it was found to be inaccurate, and recompiled and proof-read; all substantive scientific decisions were taken by the author, who assumes full responsibility for the accuracy, quality, and scientific integrity of the work.
